\section{Price/Hessian stochastic Fisher--Rao algorithms}
\label{sec:stochastic}

The deterministic maps of \Cref{sec:geometry-fr-algorithms} require Gaussian
expectations.  We now replace them by minibatch estimates built from a joint
gradient--Hessian oracle.  The sampled Hessian is not merely a convenient
representation of Price's identity: its samplewise curvature bounds preserve
the covariance bands pathwise.  This makes the global stochastic theory
unconditional.  Locally, the score noise has unbounded support, so the
appropriate statement is instead a contraction of the killed second moment,
paired with a separate finite-horizon containment estimate.

Throughout this section,
\(\alpha\Id\preceq\nabla^2V\preceq\beta\Id\),
\(\kappa=\beta/\alpha\), and \(a_n=(m_n,C_n)\).  We write
\[
  g(a):=-\E_{\N(m,C)}\nabla V,
  \qquad
  A(a):=\E_{\N(m,C)}\nabla^2V,
  \qquad
  R(a):=C^{-1}-A(a),
\]
and recall the Fisher--Rao gradient norm
\[
  \Gradnorm(a)^2
  =g(a)^\top Cg(a)
   +\frac12\norm{C^{1/2}R(a)C^{1/2}}_{\Frob}^2.
\]

\subsection{Oracle, estimators, and covariance safety}
\label{sec:stoch-oracle}

\begin{definition}[Price/Hessian oracle]
\label{def:stoch-oracle}
At a query point \(x\in\R^d\), the oracle returns the pair
\((\nabla V(x),\nabla^2V(x))\).
\end{definition}

Fix a batch size \(B\geq1\).  Conditionally on
\(\Fil_n:=\sigma(a_0,\ldots,a_n)\), draw independent
\(Z_{n,s}\sim\N(0,\Id)\) and set
\(X_{n,s}=m_n+C_n^{1/2}Z_{n,s}\).  The sticking-the-landing
score and Price/Hessian covariance estimates are
\begin{align}
  \overline b_n
  &:=\frac1B\sum_{s=1}^B
  \left[-\nabla V(X_{n,s})+C_n^{-1/2}Z_{n,s}\right],
  \label{eq:stoch-score}\\
  \overline H_n
  &:=\frac1B\sum_{s=1}^B\nabla^2V(X_{n,s}),
  \qquad
  \overline R_n:=C_n^{-1}-\overline H_n.
  \label{eq:stoch-hessian-estimator}
\end{align}
Gaussian integration by parts and Price's identity give
\begin{equation}
  \E[\overline b_n\mid\Fil_n]=g(a_n),
  \qquad
  \E[\overline R_n\mid\Fil_n]=R(a_n).
  \label{eq:stoch-unbiased}
\end{equation}
The same samples may be used in both averages.  None of the arguments below
assumes conditional independence between the score and Hessian components.
Independence is used only across the \(B\) sample pairs.

For a stepsize \(h_n>0\), the two stochastic maps are
\begin{align}
  m_{n+1}&=m_n+h_nC_n\overline b_n,
  &
  C_{n+1}
  &=C_n^{1/2}\exp\!\left(
    h_n(\Id-C_n^{1/2}\overline H_nC_n^{1/2})
  \right)C_n^{1/2},
  \label{eq:stoch-riemannian}\\
  m_{n+1}&=m_n+h_nC_n\overline b_n,
  &
  C_{n+1}
  &=(1+h_n)(C_n^{-1}+h_n\overline H_n)^{-1},
  \label{eq:stoch-kl}
\end{align}
for the Riemannian and KL variants, respectively.  Strong log-concavity
holds samplewise:
\begin{equation}
  \alpha\Id\preceq\overline H_n\preceq\beta\Id
  \qquad\text{almost surely.}
  \label{eq:stoch-sampled-band}
\end{equation}
This property would be lost for the gradient-only Price representation of
\(A(a)\), whose individual matrices need not have bounded spectrum.

We initialize the stochastic phase by the quadratic rescue from
\Cref{sec:global}.  One deterministic oracle query at \(m_{\rm in}\) gives
\begin{equation}
  H_{\rm in}=\nabla^2V(m_{\rm in}),
  \qquad
  m_0=m_{\rm in}-H_{\rm in}^{-1}\nabla V(m_{\rm in}),
  \qquad
  C_0=H_{\rm in}^{-1}.
  \label{eq:stoch-rescue}
\end{equation}
Thus \(\beta^{-1}\Id\preceq C_0\preceq\alpha^{-1}\Id\).
For a quadratic potential, \eqref{eq:stoch-rescue} is the exact optimizing
Gaussian.

\begin{proposition}[Pathwise covariance bands]
\label{prop:stoch-covariance-bands}
Initialize by \eqref{eq:stoch-rescue}.
\begin{enumerate}[label=\textup{(\roman*)}]
\item If \eqref{eq:stoch-riemannian} is run with
  \(0<h_n\leq1/\kappa\) for every \(n\), then almost surely
  \[
    \frac1{2\beta}\Id\preceq C_n\preceq\frac1\alpha\Id
    \qquad(n\geq0).
  \]
\item If \eqref{eq:stoch-kl} is run with arbitrary \(h_n>0\), then almost
  surely
  \[
    \frac1\beta\Id\preceq C_n\preceq\frac1\alpha\Id
    \qquad(n\geq0).
  \]
\end{enumerate}
No stopping, projection, clipping, or local-containment hypothesis is used.
\end{proposition}

\begin{proof}
Condition on an arbitrary realization satisfying
\eqref{eq:stoch-sampled-band}.  For the KL update, the induction hypothesis
gives
\[
  (1+h_n)\alpha\Id
  \preceq C_n^{-1}+h_n\overline H_n
  \preceq(1+h_n)\beta\Id.
\]
Inversion proves the claimed KL band for every positive stepsize.

For the Riemannian update, put
\(B_n=C_n^{1/2}\overline H_nC_n^{1/2}\).  Since
\(e^M\succeq\Id+M\) for symmetric \(M\),
\[
  C_{n+1}^{-1}
  =C_n^{-1/2}e^{h_n(B_n-\Id)}C_n^{-1/2}
  \succeq(1-h_n)C_n^{-1}+h_n\overline H_n
  \succeq\alpha\Id,
\]
where \(h_n\leq1/\kappa\leq1\).  This proves the upper covariance bound.
It also implies \(B_n\preceq\kappa\Id\), hence
\(0\preceq h_nB_n\preceq\Id\).  Convexity of the scalar exponential on
\([0,1]\), transferred spectrally, gives
\[
  e^{h_nB_n}\preceq\Id+(e-1)h_nB_n.
\]
If \(C_n^{-1}\preceq2\beta\Id\), then
\[
  C_{n+1}^{-1}
  \preceq e^{-h_n}
  \left[C_n^{-1}+(e-1)h_n\overline H_n\right]
  \preceq\beta e^{-h_n}\left[2+(e-1)h_n\right]\Id
  \preceq2\beta\Id.
\]
The last scalar inequality holds for \(0\leq h_n\leq1\).  The rescue lies
in the stronger initial band, so induction completes the proof.
\end{proof}

\subsection{Intrinsic sticking-the-landing variance}
\label{sec:stoch-variance}

For one draw \(X\sim\N(m,C)\), write
\[
  \widehat b=-\nabla V(X)+C^{-1}(X-m),
  \qquad
  \varepsilon_b=\widehat b-g(a),
  \qquad
  \varepsilon_H=\nabla^2V(X)-A(a).
\]
The intrinsic quantity controlling the additive error is the fluctuation of
the relative Hessian.

\begin{definition}[Intrinsic Hessian fluctuation]
\label{def:stoch-intrinsic-fluctuation}
For \(a=(m,C)\), define
\begin{equation}
  \Psi(a)
  :=\E\norm{
    C^{1/2}\bigl(\nabla^2V(X)-A(a)\bigr)C^{1/2}
  }_{\Frob}^2,
  \qquad X\sim\N(m,C).
  \label{eq:stoch-Psi}
\end{equation}
\end{definition}

\begin{lemma}[Intrinsic STL variance]
\label{lem:stoch-intrinsic-variance}
The single-sample score satisfies
\begin{equation}
  \E[\varepsilon_b^\top C\varepsilon_b]
  \leq
  \norm{C^{1/2}R(a)C^{1/2}}_{\Frob}^2+\Psi(a).
  \label{eq:stoch-score-variance}
\end{equation}
If \(\zeta(a;X)=(-C\varepsilon_b,C\varepsilon_HC)\) denotes the
single-sample tangent noise, then
\begin{equation}
  \E\norm{\zeta(a;X)}_{a,\FR}^2
  \leq2\Gradnorm(a)^2+\frac32\Psi(a).
  \label{eq:stoch-tangent-variance}
\end{equation}
Consequently, the conditional tangent variance of a batch average is at
most \(B^{-1}(2\Gradnorm(a_n)^2+3\Psi(a_n)/2)\).
\end{lemma}

\begin{proof}
Set \(F(x)=-\nabla V(x)+C^{-1}(x-m)\).  The Gaussian Poincar\'e inequality,
applied componentwise to \(C^{1/2}F(X)\), yields
\[
  \E[\varepsilon_b^\top C\varepsilon_b]
  \leq\E\norm{C^{1/2}(C^{-1}-\nabla^2V(X))C^{1/2}}_{\Frob}^2.
\]
Decompose the matrix on the right as
\(C^{1/2}R(a)C^{1/2}-C^{1/2}\varepsilon_HC^{1/2}\).
The cross term has zero expectation, proving
\eqref{eq:stoch-score-variance}.  The Fisher--Rao covariance block has
weight \(1/2\), so
\[
  \E\norm{\zeta}_{a,\FR}^2
  =\E[\varepsilon_b^\top C\varepsilon_b]+\frac12\Psi(a)
  \leq\norm{C^{1/2}RC^{1/2}}_{\Frob}^2+\frac32\Psi(a).
\]
The covariance residual term is at most \(2\Gradnorm(a)^2\).  Centering and
independence across the batch divide the variance by \(B\).
\end{proof}

There are two useful ways to control \(\Psi\).  The first uses only the
curvature bounds and is self-bounding.  The second uses a Lipschitz Hessian
and retains the exact Gaussian cancellation.

\begin{lemma}[Two fluctuation regimes]
\label{lem:stoch-fluctuation-regimes}
\begin{enumerate}[label=\textup{(\roman*)}]
\item If \(C\preceq\Lambda\Id\) and \(L_\Lambda=\beta\Lambda\), then
\begin{equation}
  \Psi(a)
  \leq2L_\Lambda d
    +L_\Lambda\norm{C^{1/2}R(a)C^{1/2}}_{\Frob}^2
  \leq2L_\Lambda d+2L_\Lambda\Gradnorm(a)^2.
  \label{eq:stoch-Psi-selfbound}
\end{equation}
In particular, on either rescued upper band, \(L_\Lambda=\kappa\), and
\begin{equation}
  \E[\norm{\widehat{\grad}\calE(a)}_{a,\FR}^2\mid a]
  \leq\rho_B\Gradnorm(a)^2+\frac{3\kappa d}{B},
  \qquad
  \rho_B:=1+\frac{2+3\kappa}{B}.
  \label{eq:stoch-selfbounded-second-moment}
\end{equation}
\item Suppose
\(\norm{\nabla^2V(x)-\nabla^2V(y)}_{\op}\leq
L_H\norm{x-y}\).  Then
\begin{equation}
  \Psi(a)
  \leq L_H^2\Tr(C^2)\Tr C
  \leq L_H^2\norm C_{\op}(\Tr C)^2.
  \label{eq:stoch-Psi-anisotropic}
\end{equation}
On \(C\preceq\alpha^{-1}\Id\), writing
\(\overline L_H:=L_H/\alpha^{3/2}\),
\begin{equation}
  \Psi(a)\leq\Psi_H:=d^2\overline L_H^2.
  \label{eq:stoch-Psi-Hlip}
\end{equation}
\end{enumerate}
\end{lemma}

\begin{proof}
For the first part, let
\(B_X=C^{1/2}\nabla^2V(X)C^{1/2}\) and
\(S_A=\E B_X=C^{1/2}A(a)C^{1/2}\).  Since
\(0\preceq B_X\preceq L_\Lambda\Id\),
\[
  \Psi=\E\Tr(B_X^2)-\Tr(S_A^2)
  \leq L_\Lambda\Tr S_A.
\]
The scalar inequality \(s\leq2+(s-1)^2\), applied to the eigenvalues of
\(S_A\), gives
\(\Tr S_A\leq2d+\norm{S_A-\Id}_{\Frob}^2\).  Since
\(S_A-\Id=-C^{1/2}RC^{1/2}\), this proves
\eqref{eq:stoch-Psi-selfbound}.  Combining it with
\eqref{eq:stoch-tangent-variance} gives
\eqref{eq:stoch-selfbounded-second-moment}.

For the second part, put
\(F(Z)=C^{1/2}\nabla^2V(m+C^{1/2}Z)C^{1/2}\).  Componentwise Gaussian
Poincar\'e gives
\(\Psi\leq\E\norm{D_ZF(Z)}_{\rm HS}^2\).  For a unit vector \(u\),
\[
  D_ZF(Z)[u]=C^{1/2}M_uC^{1/2},
  \qquad
  \norm{M_u}_{\op}\leq L_H(u^\top Cu)^{1/2}.
\]
In an eigenbasis of \(C\), symmetry of \(M_u\) and
\(2c_ic_j\leq c_i^2+c_j^2\) imply
\[
  \norm{C^{1/2}M_uC^{1/2}}_{\Frob}^2
  \leq\Tr(C^2M_u^2)
  \leq\norm{M_u}_{\op}^2\Tr(C^2).
\]
Summing over an orthonormal basis of directions \(u\) proves the first
inequality in \eqref{eq:stoch-Psi-anisotropic}.  The remaining inequalities
follow from
\(\Tr(C^2)\leq\norm C_{\op}\Tr C\) and the rescued upper band.
\end{proof}

The anisotropic form in \eqref{eq:stoch-Psi-anisotropic} can be much smaller
than \(d^2\overline L_H^2\) when the covariance spectrum decays.  The
worst-case \(d^2\) bound is nevertheless the uniform consequence of the
band.  For a quadratic target, \(L_H=0\) and, more strongly, the integrand in
\eqref{eq:stoch-Psi} is identically zero.

\subsection{Global stochastic convergence}
\label{sec:stoch-global}

On the Riemannian band of \Cref{prop:stoch-covariance-bands}, the global
gradient-domination estimate from \Cref{sec:global} reads
\(\Gradnorm(a_n)^2\geq\Delta_n/(2\kappa)\).  On the tighter KL band it reads
\(\Gradnorm(a_n)^2\geq\Delta_n/\kappa\), where
\(\Delta_n=\calE(a_n)-\calE(a_\star)\).

\begin{lemma}[Intrinsic one-step inequalities]
\label{lem:stoch-global-onestep}
Conditionally on \(\Fil_n\), the Riemannian update with
\(0<h\leq1/\kappa\) satisfies
\begin{equation}
  \E[\Delta_{n+1}\mid\Fil_n]
  \leq\Delta_n
  -h\left[1-\kappa h\left(1+\frac2B\right)\right]
     \Gradnorm(a_n)^2
  +\frac{3\kappa h^2}{2B}\Psi(a_n).
  \label{eq:stoch-riemannian-onestep}
\end{equation}
The KL update with \(0<h\leq1/(8\kappa)\) satisfies
\begin{equation}
  \E[\Delta_{n+1}\mid\Fil_n]
  \leq\Delta_n-\frac h2\Gradnorm(a_n)^2
  +\frac{7\kappa h^2}{4B}\Psi(a_n).
  \label{eq:stoch-kl-onestep}
\end{equation}
Both estimates remain valid when the score and Hessian use the same batch.
\end{lemma}

\begin{proof}
The retraction smoothness estimate from the deterministic analysis gives,
pathwise on \(C_n\preceq\alpha^{-1}\Id\),
\[
  \calE(a_{n+1})
  \leq\calE(a_n)
  -h\inner{\grad\calE(a_n)}{\widehat{\grad}\calE(a_n)}_{a_n}
  +\kappa h^2\norm{\widehat{\grad}\calE(a_n)}_{a_n,\FR}^2.
\]
Taking conditional expectation, using unbiasedness and
\Cref{lem:stoch-intrinsic-variance}, proves
\eqref{eq:stoch-riemannian-onestep}.

For the KL map, suppress the time index and set
\[
  S=C^{1/2}R(a)C^{1/2},
  \quad
  D=C^{1/2}(\overline H-A(a))C^{1/2},
  \quad
  P=C^{1/2}\overline HC^{1/2},
  \quad
  \overline S=S-D.
\]
The covariance update has the exponential-retraction representation
\begin{equation}
  C^+=C^{1/2}e^{W(P)}C^{1/2},
  \qquad
  W(P)=\log(1+h)\Id-\log(\Id+hP).
  \label{eq:stoch-kl-logarithmic-tangent}
\end{equation}
For every eigenvalue \(p\in[1/\kappa,\kappa]\), the integral formula
\(W(p)=\int_p^1h(1+hr)^{-1}\dd r\) gives
\begin{equation}
  \norm{W(P)}_{\Frob}\leq h\norm{\overline S}_{\Frob},
  \qquad
  \norm{W(P)-h\overline S}_{\Frob}
  \leq\kappa h^2\norm{\overline S}_{\Frob}.
  \label{eq:stoch-kl-tangent-bounds}
\end{equation}
Write
\(M=g^\top Cg\), \(V=\norm S_{\Frob}^2\), and
\(\sigma^2=\Psi(a)/B\).  Then
\(\E\norm D_{\Frob}^2=\sigma^2\) and
\eqref{eq:stoch-score-variance} gives
\(\E[\overline b^\top C\overline b]\leq M+V/B+\sigma^2\).
Applying the same retraction smoothness bound to the logarithmic tangent in
\eqref{eq:stoch-kl-logarithmic-tangent}, and then using
\eqref{eq:stoch-kl-tangent-bounds}, yields
\[
\begin{split}
  \E[\Delta^+\mid\Fil]
  \leq{}&\Delta-h(1-\kappa h)M
  -h\left[\frac12-\kappa h\left(1+\frac1B\right)\right]V
  +\frac74\kappa h^2\sigma^2.
\end{split}
\]
This calculation uses only \(\E D=0\) and the two marginal second moments;
there is no score--Hessian independence step.  If
\(h\leq1/(8\kappa)\), the coefficients of \(M\) and \(V\) are at least
\(h/2\) and \(h/4\), respectively.  Since
\(\Gradnorm^2=M+V/2\), \eqref{eq:stoch-kl-onestep} follows.
\end{proof}

The next result displays the two fluctuation regimes side by side.  Every
constant-step conclusion is a geometric approach to the stated residual
upper bound.  It is not a claim of convergence to arbitrary accuracy, nor a
matching lower bound on the algorithm's stationary error.

\begin{theorem}[Constant-step global rates and residual floors]
\label{thm:stoch-global-constant-step}
Initialize by \eqref{eq:stoch-rescue} and let \(B\geq1\).
\begin{enumerate}[label=\textup{(\roman*)}]
\item Under the Hessian-Lipschitz assumption, the Riemannian method with
\(h=[2\kappa(1+2/B)]^{-1}\) satisfies
\begin{equation}
  \E\Delta_N
  \leq
  \left(1-\frac1{8\kappa^2(1+2/B)}\right)^N\Delta_0
  +\frac{3\kappa}{B+2}\,d^2\overline L_H^2.
  \label{eq:stoch-global-riemannian-Hlip}
\end{equation}
\item Under the same assumption, the KL method with
\(h=1/(8\kappa)\) satisfies
\begin{equation}
  \E\Delta_N
  \leq
  \left(1-\frac1{16\kappa^2}\right)^N\Delta_0
  +\frac{7\kappa}{16B}\,d^2\overline L_H^2.
  \label{eq:stoch-global-kl-Hlip}
\end{equation}
\item Without Hessian-Lipschitzness, set
\(\rho_B=1+(2+3\kappa)/B\).  The Riemannian method with
\(h=(2\kappa\rho_B)^{-1}\) satisfies
\begin{equation}
  \E\Delta_N
  \leq
  \left(1-\frac1{8\kappa^2\rho_B}\right)^N\Delta_0
  +\frac{6\kappa^2d}{B+2+3\kappa}.
  \label{eq:stoch-global-riemannian-selfbound}
\end{equation}
\item Without Hessian-Lipschitzness, the KL method with
\(h=\min\{1/(8\kappa),B/(14\kappa^2)\}\) satisfies
\begin{equation}
  \E\Delta_N
  \leq
  \left(1-\frac h{4\kappa}\right)^N\Delta_0
  +\frac{14\kappa^3hd}{B}.
  \label{eq:stoch-global-kl-selfbound}
\end{equation}
If \(B\leq7\kappa/4\), the last floor is exactly \(\kappa d\) and the
contraction factor is \(1-B/(56\kappa^3)\).
\end{enumerate}
\end{theorem}

\begin{proof}
For (i), the bracket in
\eqref{eq:stoch-riemannian-onestep} equals \(1/2\).  Use
\(\Psi\leq\Psi_H=d^2\overline L_H^2\) and
\(\Gradnorm^2\geq\Delta/(2\kappa)\), then sum the geometric recursion.
Its additive increment is \(3\kappa h^2\Psi_H/(2B)\), and division by the
drift \(h/(4\kappa)\) gives the floor in
\eqref{eq:stoch-global-riemannian-Hlip}.  Part (ii) follows in the same way
from \eqref{eq:stoch-kl-onestep} and
\(\Gradnorm^2\geq\Delta/\kappa\); the drift is
\(h/(2\kappa)\).

For (iii), substitute
\(\Psi\leq2\kappa d+2\kappa\Gradnorm^2\) into
\eqref{eq:stoch-riemannian-onestep}.  This gives
\[
  \E[\Delta_{n+1}\mid\Fil_n]
  \leq\Delta_n-h(1-\kappa h\rho_B)\Gradnorm(a_n)^2
  +\frac{3\kappa^2h^2d}{B}.
\]
At \(h=(2\kappa\rho_B)^{-1}\), the first bracket is \(1/2\).  Applying
the Riemannian-band gradient domination and unrolling gives
\eqref{eq:stoch-global-riemannian-selfbound}.  For (iv), substitution into
\eqref{eq:stoch-kl-onestep} gives
\[
  \E[\Delta_{n+1}\mid\Fil_n]
  \leq\Delta_n
  -\frac h2\left(1-\frac{7\kappa^2h}{B}\right)\Gradnorm(a_n)^2
  +\frac{7\kappa^2h^2d}{2B}.
\]
The chosen stepsize makes the bracket at least \(1/2\).  The KL-band
gradient domination and a geometric sum prove
\eqref{eq:stoch-global-kl-selfbound}.
\end{proof}

For \(B=1\), the Hessian-Lipschitz bounds require
\(\widetilde O(\kappa^2)\) rounds to reach an accuracy above their displayed
floors.  Under curvature bounds alone, both single-sample bounds require
\(\widetilde O(\kappa^3)\) rounds.  Taking
\(B=\Theta(\kappa)\) reduces the latter round complexity to
\(\widetilde O(\kappa^2)\), but the number of Price/Hessian oracle pairs
remains \(\widetilde O(\kappa^3)\).  The one-time rescue adds one oracle
pair.

Constant steps are not necessary if convergence all the way to zero is
required in the Hessian-Lipschitz regime.

\begin{theorem}[Floor-free decreasing steps]
\label{thm:stoch-global-decreasing}
Under the Hessian-Lipschitz assumption, initialize by
\eqref{eq:stoch-rescue} and run either method with
\begin{equation}
  h_n=\frac{8\kappa}{n+n_0},
  \qquad n_0\geq64\kappa^2.
  \label{eq:stoch-decreasing-schedule}
\end{equation}
Then, for all \(N\geq0\),
\begin{equation}
  \E\Delta_N\leq\frac K{N+n_0},
  \qquad
  K:=\max\left\{
    \frac{112\kappa^3d^2\overline L_H^2}{B},
    n_0\Delta_0
  \right\}.
  \label{eq:stoch-decreasing-rate}
\end{equation}
Thus \eqref{eq:stoch-decreasing-rate} is a genuine \(O(N^{-1})\) decay to
zero.  The algorithm does not need to know \(L_H\); that constant appears
only in the certificate.
\end{theorem}

\begin{proof}
The schedule satisfies \(h_n\leq1/(8\kappa)\).  The KL band persists for
all positive steps.  The proof of
\Cref{prop:stoch-covariance-bands} also works step by step for the
Riemannian schedule because each \(h_n\leq1/\kappa\).  Hence
\eqref{eq:stoch-riemannian-onestep} and
\eqref{eq:stoch-kl-onestep} apply at every iteration.  In both cases, using
the weaker of the two band-dependent drift constants and
\(\Psi\leq\Psi_H\),
\begin{equation}
  \E\Delta_{n+1}
  \leq
  \left(1-\frac{h_n}{4\kappa}\right)\E\Delta_n
  +\frac{7\kappa h_n^2}{4B}\Psi_H.
  \label{eq:stoch-decreasing-recursion}
\end{equation}
With \(t=n+n_0\), this is
\[
  \E\Delta_{n+1}
  \leq\left(1-\frac2t\right)\E\Delta_n+\frac S{t^2},
  \qquad
  S=\frac{112\kappa^3\Psi_H}{B}.
\]
If \(K=\max\{S,n_0\Delta_0\}\), induction gives
\(\E\Delta_n\leq K/(n+n_0)\): assuming the claim at \(n\),
\[
  \E\Delta_{n+1}
  \leq \frac Kt-\frac{2K-S}{t^2}
  \leq \frac{K(t-1)}{t^2}
  \leq \frac K{t+1}.
\]
This proves \eqref{eq:stoch-decreasing-rate}.
\end{proof}

\subsection{Local killed-process contraction and containment}
\label{sec:stoch-local}

We now work in the optimizer-whitened coordinates of \Cref{sec:local}, so
\(a_\star=(0,\Id)\).  The two equilibrium norms are
\begin{equation}
  \norm{(m,C)-a_\star}_{\star,\Riem}^2
  :=\norm m^2+\frac12\norm{C-\Id}_{\Frob}^2,
  \qquad
  \norm{(m,C)-a_\star}_{\star,\KL}^2
  :=\norm m^2+\frac{1+h}{2}\norm{C-\Id}_{\Frob}^2.
  \label{eq:stoch-local-norms}
\end{equation}
If the original Hessian has Lipschitz modulus \(L_H\), then the whitened
potential has modulus
\begin{equation}
  L_{H,\star}
  \leq L_H\norm{C_\star}_{\op}^{3/2}
  \leq \frac{L_H}{\alpha^{3/2}}
  =\overline L_H.
  \label{eq:stoch-whitened-Hlip}
\end{equation}

For a fixed method and stepsize, call a local energy level
\(\delta_\bullet\) stochastically admissible if the deterministic
second-derivative criterion of \Cref{sec:local-discrete} holds and, on
\(\mathcal U_{\delta_\bullet}\),
\begin{equation}
  \frac12\Id\preceq C\preceq\frac32\Id,
  \qquad
  \norm{a-a_\star}_{\star,\bullet}\leq1.
  \label{eq:stoch-local-admissible}
\end{equation}
Equivalently, one may require
\(K_{\mathrm{disc},\bullet}(\delta_\bullet)
\mathfrak b_\bullet(\delta_\bullet)\sqrt{\delta_\bullet}
\leq\gamma_\bullet\), with
\(\gamma_{\Riem}=1/(4+\Gamma)\) and
\(\gamma_{\KL}=1/(4+2h+\Gamma)\), together with
\eqref{eq:stoch-local-admissible}.  Such levels exist for every fixed
admissible \(h\) by shrinking \(\delta_\bullet\).  As in
\Cref{rem:local-discrete-explicitness}, this is a constructive existence
statement through \(K_{\mathrm{disc},\bullet}\), not a closed-form radius in the target
parameters.

The local noise consists of the Euclidean score error
\(\varepsilon_b\) and the whitened Hessian error \(\varepsilon_H\).

\begin{lemma}[Local variance decomposition]
\label{lem:stoch-local-variance}
If \(\Id/2\preceq C\preceq3\Id/2\), then
\begin{equation}
  \E\left[\norm{\varepsilon_b}^2
    +\frac12\norm{\varepsilon_H}_{\Frob}^2\right]
  \leq V_0+V_1\norm{a-a_\star}_{\star,\Riem}^2,
  \qquad
  V_0:=3d^2L_{H,\star}^2,
  \quad
  V_1:=24+6dL_{H,\star}^2.
  \label{eq:stoch-local-variance}
\end{equation}
For a batch average, the right-hand side is divided by \(B\).  At
\(a_\star\), the sharper bound is
\[
  \E\left[\norm{\varepsilon_b}^2
    +\frac12\norm{\varepsilon_H}_{\Frob}^2\right]
  \leq\frac32d^2L_{H,\star}^2.
\]
\end{lemma}

\begin{proof}
Gaussian Poincar\'e applied to the Hessian entries gives
\(\E\norm{\varepsilon_H}_{\Frob}^2\leq(3/2)d^2L_{H,\star}^2\).
Applying it to
\(F(x)=-\nabla\widehat V(x)+C^{-1}(x-m)\) gives
\[
  \E\norm{\varepsilon_b}^2
  \leq\frac32\left(
    \norm{C^{-1}-A(a)}_{\Frob}^2
    +\E\norm{\varepsilon_H}_{\Frob}^2
  \right).
\]
On the stated band,
\(\norm{C^{-1}-\Id}_{\Frob}\leq2\norm{C-\Id}_{\Frob}\).
Moreover, coupling \(m+C^{1/2}Z\) with \(Z\) and using
\(A(a_\star)=\Id\) yields
\[
  \norm{A(a)-\Id}_{\Frob}^2
  \leq dL_{H,\star}^2
  \left(\norm m^2+\norm{C^{1/2}-\Id}_{\Frob}^2\right).
\]
Since
\(\norm{C^{1/2}-\Id}_{\Frob}\leq\norm{C-\Id}_{\Frob}\), collecting the
terms gives \eqref{eq:stoch-local-variance}.  At equilibrium the residual
\(C^{-1}-A(a)\) vanishes and \(\norm C_{\op}=1\), which gives the sharper
constant.  Batch averaging divides centered second moments by \(B\).
\end{proof}

Let \(T_{\bullet,h}\) denote the corresponding deterministic local map.
The following perturbation estimate records the constants needed below.

\begin{lemma}[One-step stochastic perturbation]
\label{lem:stoch-local-perturbation}
Suppose \eqref{eq:stoch-local-admissible} holds and
\(\norm{T_{\bullet,h}(a)-a_\star}_{\star,\bullet}\leq1\), with
\(0<h\leq1\).  One stochastic step satisfies
\begin{align}
  \E[\norm{a^+-a_\star}_{\star,\Riem}^2\mid a]
  &\leq
  \norm{T_{\Riem,h}(a)-a_\star}_{\star,\Riem}^2
  +\frac{64h^2}{B}\left(V_0+V_1\norm{a-a_\star}_{\star,\Riem}^2\right),
  \label{eq:stoch-local-perturb-riemannian}\\
  \E[\norm{a^+-a_\star}_{\star,\KL}^2\mid a]
  &\leq
  \norm{T_{\KL,h}(a)-a_\star}_{\star,\KL}^2
  +\frac{68h^2}{B}\left(V_0+V_1\norm{a-a_\star}_{\star,\Riem}^2\right).
  \label{eq:stoch-local-perturb-kl}
\end{align}
\end{lemma}

\begin{proof}
For the Riemannian covariance update, put
\(S=\Id-C^{1/2}A(a)C^{1/2}\) and replace it in the stochastic map by
\(\widehat S=S-C^{1/2}\overline\varepsilon_HC^{1/2}\).  For
\(S,\widehat S\preceq\Id\) and \(h\leq1\), the Fr\'echet derivative and its
integral remainder give
\begin{align*}
  \norm{e^{h\widehat S}-e^{hS}}_{\Frob}
  &\leq eh\norm{\widehat S-S}_{\Frob},\\
  \norm{e^{h\widehat S}-e^{hS}
       -D\exp(hS)[h(\widehat S-S)]}_{\Frob}
  &\leq\frac e2h^2
  \norm{\widehat S-S}_{\op}\norm{\widehat S-S}_{\Frob}.
\end{align*}
The linear term is centered.  Expanding the squared norm around the
deterministic update, the second-moment and bias coefficients are bounded by
\(81e^2/16\) and \(27e/(4\sqrt2)\), whose sum is less than \(64\).

For KL, write
\(P=C^{-1}+hA(a)\) and
\(\widehat P=P+h\overline\varepsilon_H\).  The identities
\[
  \widehat P^{-1}-P^{-1}
  =-P^{-1}(h\overline\varepsilon_H)\widehat P^{-1}
\]
and
\[
  \widehat P^{-1}-P^{-1}
  =-P^{-1}(h\overline\varepsilon_H)P^{-1}
   +P^{-1}(h\overline\varepsilon_H)P^{-1}
    (h\overline\varepsilon_H)\widehat P^{-1}
\]
control the second moment and the second-order bias.  On the covariance
band, their aggregate coefficient is
\(81/2+27=67.5<68\).  The score increment is linear and centered in both
maps.  Applying \Cref{lem:stoch-local-variance} completes the proof; in the
KL case one also uses
\(\norm{a-a_\star}_{\star,\Riem}\leq
\norm{a-a_\star}_{\star,\KL}\).
\end{proof}

Set
\begin{equation}
\begin{array}{c|c|c|c}
  \bullet&A_\bullet&c_\bullet&
  \eta_\bullet\\ \hline
  \Riem&4+\Gamma&64&
  192h^2d^2L_{H,\star}^2/B\\[1mm]
  \KL&4+2h+\Gamma&68&
  204h^2d^2L_{H,\star}^2/B
\end{array}
\qquad
\mu_\bullet:=\frac1{2A_\bullet}.
\label{eq:stoch-local-constants}
\end{equation}

\begin{theorem}[Local non-exit mean-square contraction]
\label{thm:stoch-local-killed}
Fix \(\bullet\in\{\Riem,\KL\}\) and suppose
\begin{equation}
  0<h\leq
  \min\left\{
    \frac2{4+\Gamma},
    \frac{B}{4c_\bullet A_\bullet V_1}
  \right\}.
  \label{eq:stoch-local-stepsize}
\end{equation}
Let \(\delta_\bullet\) be stochastically admissible for this \(h\), start
in \(\mathcal U_{\delta_\bullet}\), and define
\(\tau_\bullet=\inf\{n\geq0:a_n\notin
\mathcal U_{\delta_\bullet}\}\).  Then, for every \(N\geq0\),
\begin{equation}
  \E\left[
    \norm{a_N-a_\star}_{\star,\bullet}^2
    \one_{\{\tau_\bullet>N\}}
  \right]
  \leq
  (1-\mu_\bullet h)^N
  \norm{a_0-a_\star}_{\star,\bullet}^2
  +\frac{\eta_\bullet}{\mu_\bullet h}.
  \label{eq:stoch-local-killed}
\end{equation}
In particular, the two residual terms are
\begin{align}
  \frac{\eta_{\Riem}}{\mu_{\Riem}h}
  &=\frac{384h(4+\Gamma)d^2L_{H,\star}^2}{B},
  \label{eq:stoch-local-floor-riemannian}\\
  \frac{\eta_{\KL}}{\mu_{\KL}h}
  &=\frac{408h(4+2h+\Gamma)d^2L_{H,\star}^2}{B}.
  \label{eq:stoch-local-floor-kl}
\end{align}
The left-hand side of \eqref{eq:stoch-local-killed} is a killed, or
non-exit, second moment.  It is not
\(\E\norm{a_{N\wedge\tau_\bullet}-a_\star}_{\star,\bullet}^2\).
\end{theorem}

\begin{proof}
Write \(r_n=\norm{a_n-a_\star}_{\star,\bullet}\).  On
\(\{n<\tau_\bullet\}\), the deterministic second-derivative criterion and
the linearized bounds of \Cref{cor:local-linear-maps} give
\[
  \norm{T_{\bullet,h}(a_n)-a_\star}_{\star,\bullet}
  \leq(1-\mu_\bullet h)r_n.
\]
Apply \Cref{lem:stoch-local-perturbation}.  Since
\(\mu_\bullet h\leq1/2\),
\((1-\mu_\bullet h)^2\leq1-3\mu_\bullet h/2\).  The second restriction in
\eqref{eq:stoch-local-stepsize} gives
\(c_\bullet h^2V_1/B\leq\mu_\bullet h/2\).  Therefore
\begin{equation}
  \E[r_{n+1}^2\mid\Fil_n]
  \leq(1-\mu_\bullet h)r_n^2+\eta_\bullet
  \qquad\text{on }\{n<\tau_\bullet\}.
  \label{eq:stoch-local-one-step-recursion}
\end{equation}
Let
\(M_n=\E[r_n^2\one_{\{\tau_\bullet>n\}}]\).  Because
\(\{\tau_\bullet>n+1\}\subseteq\{\tau_\bullet>n\}\),
\(M_{n+1}\leq(1-\mu_\bullet h)M_n+\eta_\bullet\).
Summing this geometric recursion proves
\eqref{eq:stoch-local-killed}; substituting
\eqref{eq:stoch-local-constants} gives
\eqref{eq:stoch-local-floor-riemannian}--\eqref{eq:stoch-local-floor-kl}.
\end{proof}

The missing mass in \eqref{eq:stoch-local-killed} must be controlled
separately.  The next proposition does so without a union bound over all
possible exit times.

\begin{proposition}[Finite-horizon containment]
\label{prop:stoch-local-containment}
Choose \(r\leq1/(2\sqrt2)\) such that the closed
\(\norm{\cdot}_{\star,\bullet}\)-ball of radius \(r\) about \(a_\star\) is
contained in \(\mathcal U_{\delta_\bullet}\), and set
\[
  \tau_r:=\inf\{n\geq0:
    \norm{a_n-a_\star}_{\star,\bullet}>r\}.
\]
Under the assumptions of \Cref{thm:stoch-local-killed}, if
\(r_0:=\norm{a_0-a_\star}_{\star,\bullet}\leq r\), then
\begin{equation}
  \Prob(\tau_r\leq N)
  \leq\frac{r_0^2+N\eta_\bullet}{r^2}
  \qquad(N\geq1).
  \label{eq:stoch-local-exit-probability}
\end{equation}
Consequently, \(\Prob(\tau_r\leq N)\leq\varepsilon_{\rm exit}\) follows
from the two separate conditions
\begin{equation}
  r_0^2\leq\frac{\varepsilon_{\rm exit}}2r^2,
  \qquad
  B\geq
  \frac{2c'_\bullet Nh^2d^2L_{H,\star}^2}
       {\varepsilon_{\rm exit}r^2},
  \quad
  c'_{\Riem}=192,
  \quad c'_{\KL}=204.
  \label{eq:stoch-local-containment-conditions}
\end{equation}
Thus increasing the batch cannot compensate for an initial point near the
outer boundary.  Deep entry is necessary in this argument, and the batch
needed for survival grows linearly with the prescribed horizon \(N\).
\end{proposition}

\begin{proof}
Freeze the radial process at the exit time:
\(R_n=r_{n\wedge\tau_r}\).  On \(\{\tau_r>n\}\),
\eqref{eq:stoch-local-one-step-recursion} and
\(1-\mu_\bullet h\leq1\) give
\(\E[R_{n+1}^2\mid\Fil_n]\leq R_n^2+\eta_\bullet\).  After exit, the
process is constant, so the same inequality remains true.  Hence
\(\E R_N^2\leq r_0^2+N\eta_\bullet\).  On
\(\{\tau_r\leq N\}\), one has \(R_N>r\); Markov's inequality proves
\eqref{eq:stoch-local-exit-probability}.  Substitution of
\eqref{eq:stoch-local-constants} gives
\eqref{eq:stoch-local-containment-conditions}.
\end{proof}

The certified local round complexity follows directly.  Up to numerical
constants, the largest permitted stepsize is
\[
  h_\bullet^\star
  =\Theta\!\left(
    \frac1{4+\Gamma}\min\left\{1,\frac B{V_1}\right\}
  \right),
\]
where the KL choice may use \(4+\Gamma+1\) in place of its mildly implicit
\(A_{\KL}\).  Therefore
\begin{equation}
  N_{\rm loc}
  =O\!\left(
    (4+\Gamma)^2
    \max\left\{1,\frac{V_1}{B}\right\}
    \log\frac{r_0^2}{\varepsilon}
  \right).
  \label{eq:stoch-local-complexity}
\end{equation}
In the large-batch branch and with the universal spectral certificate
\(\Gamma=O(1+\sqrt d\log\kappastar)\), this becomes
\[
  O\!\left((1+d\log^2\kappastar)
  \log\frac{r_0^2}{\varepsilon}\right).
\]
The dependence is quadratic in \(4+\Gamma\), not linear.  In particular,
the present theorem does not prove an unqualified logarithmic dependence on
\(\kappastar\).  Moreover, \eqref{eq:stoch-local-complexity} is a round
count for contraction of the killed moment.  Enforcing survival over those
rounds may increase the batch according to
\eqref{eq:stoch-local-containment-conditions}, so it is not by itself a
high-probability oracle complexity.

The global and local statements can nevertheless be composed without
hiding these qualifications.

\begin{corollary}[Stochastic global-to-local composition]
\label{cor:stoch-global-to-local}
Let an inner level \(\delta_{\rm in}\) and an outer radius \(r\) satisfy
\[
  \mathfrak b_\bullet(\delta_{\rm in})^2\delta_{\rm in}
  \leq\frac{\varepsilon_{\rm exit}}2r^2,
\]
with the radius-\(r\) ball contained in a stochastically admissible outer
level.  Suppose a global phase produces
\begin{equation}
  \E\Delta_{N_1}
  \leq\varepsilon_{\rm ent}\delta_{\rm in}.
  \label{eq:stoch-global-to-local-entry}
\end{equation}
For a constant-step global method, its residual floor in
\Cref{thm:stoch-global-constant-step} must be below the right-hand side of
\eqref{eq:stoch-global-to-local-entry}; the decreasing schedule of
\Cref{thm:stoch-global-decreasing} has no such obstruction.  Run the local
method for \(N_2\) further steps with the batch required by
\eqref{eq:stoch-local-containment-conditions}.  Then, with
\(E_{\rm in}=\{a_{N_1}\in\mathcal U_{\delta_{\rm in}}\}\) and the local
exit time shifted to \(N_1\),
\begin{equation}
  \Prob(E_{\rm in}\text{ and no exit through }N_2)
  \geq1-\varepsilon_{\rm ent}-\varepsilon_{\rm exit}.
  \label{eq:stoch-global-to-local-survival}
\end{equation}
The entry-restricted killed second moment satisfies
\begin{equation}
  \E\!\left[
    \norm{a_{N_1+N_2}-a_\star}_{\star,\bullet}^2;
    E_{\rm in},\ \text{no exit through }N_2
  \right]
  \leq(1-\mu_\bullet h)^{N_2}
  \mathfrak b_\bullet(\delta_{\rm in})^2\delta_{\rm in}
  +\frac{\eta_\bullet}{\mu_\bullet h}.
  \label{eq:stoch-global-to-local-moment}
\end{equation}
\end{corollary}

\begin{proof}
Markov's inequality and \eqref{eq:stoch-global-to-local-entry} give
\(\Prob(E_{\rm in}^c)\leq\varepsilon_{\rm ent}\).  On
\(E_{\rm in}\), the norm-conversion estimate of \Cref{sec:local} gives
the deep-entry condition in
\eqref{eq:stoch-local-containment-conditions}; that proposition bounds the
conditional exit mass by \(\varepsilon_{\rm exit}\).  A union bound proves
\eqref{eq:stoch-global-to-local-survival}.  Iterating
\eqref{eq:stoch-local-one-step-recursion} after time \(N_1\), with the
indicator of \(E_{\rm in}\) retained, gives
\eqref{eq:stoch-global-to-local-moment}.
\end{proof}

\subsection{Gaussian calibration}
\label{sec:stoch-gaussian}

The preceding floors measure nonconstant Hessian fluctuations, not Monte
Carlo noise in the abstract.  Gaussian targets make the distinction exact.

\begin{proposition}[Gaussian target calibration]
\label{prop:stoch-gaussian-calibration}
Let
\(V(x)=\tfrac12(x-\mu)^\top H(x-\mu)+c\), \(H\succ0\).
\begin{enumerate}[label=\textup{(\roman*)}]
\item The Hessian is constant, hence \(\Psi(a)=0\) for every \(a\).  The
additive terms in the intrinsic one-step inequalities
\eqref{eq:stoch-riemannian-onestep} and
\eqref{eq:stoch-kl-onestep} vanish exactly.  The nonzero floors displayed by
the coarse self-bounding bounds are artifacts of substituting
\eqref{eq:stoch-Psi-selfbound}, not actual Gaussian floors.
\item In optimizer-whitened coordinates, the single-sample score noise is
\begin{equation}
  \varepsilon_b=(C^{-1/2}-C^{1/2})Z.
  \label{eq:stoch-gaussian-score-noise}
\end{equation}
It is nonzero and has unbounded support whenever \(C\neq\Id\).  Thus zero
additive floor does not imply pathwise deterministic evolution from an
arbitrary covariance.
\item If \(C_0=C_\star\), then \(C_n=C_\star\) for both methods, all sampled
noise vanishes, and in whitened coordinates
\begin{equation}
  m_{n+1}=(1-h_n)m_n
  \label{eq:stoch-gaussian-matched-covariance}
\end{equation}
pathwise.  If the quadratic rescue is used, it returns
\((m_0,C_0)=(\mu,H^{-1})=a_\star\) after its single query, so there is no
iterative phase.
\item Locally, \(L_{H,\star}=0\), hence \(V_0=\eta_\bullet=0\).  The killed
contraction has no additive term and the containment bound loses its
\(N\)-growing term.  For a generic unmatched covariance, however, these
remain mean-square and probability statements rather than pathwise
containment because of \eqref{eq:stoch-gaussian-score-noise}.
\end{enumerate}
\end{proposition}

\begin{proof}
Constancy of \(\nabla^2V=H\) makes the integrand in
\eqref{eq:stoch-Psi} zero.  After whitening, the target is standard
Gaussian and \(X=m+C^{1/2}Z\).  Since \(g(a)=-m\),
\[
  \widehat b-g(a)
  =-C^{1/2}Z+C^{-1/2}Z,
\]
which proves \eqref{eq:stoch-gaussian-score-noise}.  When \(C=\Id\), this
noise vanishes and \(\overline H=\Id\) pathwise.  Both covariance maps then
fix \(\Id\), while their common mean update is
\eqref{eq:stoch-gaussian-matched-covariance}.  Finally, direct substitution
into \eqref{eq:stoch-rescue} gives \(m_0=\mu\), \(C_0=H^{-1}\), and
substitution of \(L_{H,\star}=0\) into the local constants proves the last
claim.
\end{proof}
